\chapter*{使用说明}
\addcontentsline{toc}{chapter}{使用说明}
\markboth{使用说明}{使用说明}
\begin{enumerate}
\item \key{定位}：本册只收“考场上要查的东西”——定义、假设、模型、估计量公式、方差与标准误、检验统计量及其分布与自由度、成立条件、方法对比表、关键推导骨架、易错点、英文一句话表述。不收例题与练习；第 17 章“2025 年真题精讲”保留每一小问的完整推导与英文标准答案。
\item \key{与教材节号一致}：本册的章、节编号与《高级计量经济学 复习教材（新版）》完全对应（本册 §10.6 = 教材 §10.6）。需要完整推导、例题或更多解释时，直接翻到教材同一节。
\item \key{怎么查}：目录细到小知识点；页眉显示章名（左页）与节名（右页）。英文术语查附录“中英术语速查”；临界值查附录 A（表前有“统计量—分布—查哪张表”总览）；英文作答句式查附录 B。
\item \key{标记}：推导右侧灰色〔……〕为该步理由；\tagext\ 为拓展内容；\tagexam\ 为与真题直接相关的内容；“易错”条目列出最常见的错误写法。
\item \key{记号}：$X\T$ 为转置；$\one_T$ 为 $T$ 维全 1 向量；$P_X=X(X\T X)\inv X\T$，$M_X=I-P_X$；$\pto$、$\dto$ 为依概率、依分布收敛；$\E$ 为期望。
\end{enumerate}
\begin{center}
\begin{tabularx}{\linewidth}{>{\bfseries}p{3.3cm}L}
\toprule
题目问的是 & 先查 \\
\midrule
推导估计量的公式 & 对应章“估计量”小节的推导骨架（第 3、8、10–12 章） \\
证明一致性/渐近正态 & 第 5 章“分母/分子”套路；GMM 见第 8 章；面板见第 10–12 章 \\
写标准误/检验统计量 & 第 5 章三种标准误；第 6 章 $t$/$F$/Wald；第 9 章 IV 三大检验；第 11 章 Hausman \\
选择估计方法 & 各章末“本章小结”的方法对照表；第 15–16 章处理效应方法总表 \\
DiD / TWFE / 事件研究 & 第 15 章；真题第 17 章 \\
英文作答 & 各节 English 表述；附录 B 作答结构与句式 \\
\bottomrule
\end{tabularx}
\end{center}
